SEE provides an interactive plot alongside a sidebar that guides users from first contact---via
a tutorial---through placing voters and candidates, running elections under various voting rules,
and optionally grounding the simulation in real political survey data.

\subsection{Tutorial}%
\label{sub:Tutorial}
The tutorial serves to orient the user, introduce the voting rules,
and set user expectations for the allowable interaction patterns in the application.

It begins by defining the spatial model, contextualizing the axes
as representing a policy or ideology space, and emphasizing the
key assumption: only the distance between voters and candidates matters
in determining how voters will fill out their ballots.

Next, it explains the application and the primary ways of using it
via a series of examples that introduce Plurality, Borda, IRV, and STV.
At each step, the tutorial shows an example plot
and asks the user to make a prediction of which candidate will win.

As users walk through the examples in the tutorial tab, they will learn about the "spoiler effect" in 
plurality voting systems, and how Instant-Runoff voting (IRV) avoids this effect. 
Next, a lesson on how Borda scores are calculated and how those scores result in electing 
centrist "consensus" candidates. Lastly, within the the multi-winner setting, users discover that
 Single Transferable Vote (STV) can elect a more representative set of candidates than Bloc Plurality.

\subsection{Input}%
\label{sub:Input}
On the plot, voters are blue dots, and candidates are orange stars.
There are four primary methods for placing voters and candidates on the plot.
\begin{itemize}
  \item \textbf{Manual Input.}
    In the upper left, users can click either the "+candidate" or "+voter" button, or engage these menus
    by typing 'c' or 'v' on the keyboard, respectively. Then users can click anywhere in the 
    plot to add a voter or candidate point at that location. To exit, users right click or press 'esc' to exit.
    Note that within the "Survey" tab, the add voter menu is disabled, as each voter point represents a real survey respondent.

  \item \textbf{Synthetic Sampling.}
    Within the "Synthetic" tab, the user can choose one of the available distributions.
    The distribution center is set by dragging it on the plot, or through the input boxes.
    The distribution parameters adjust with slider(s) or through inputting values to the right of the slider(s).
    Then users then selects the type of point, voters or candidates, to be randomly sampled from the specified distribution, 
    and inputs the desired count. Clicking "add data" adds these points to the plot.
    %
    As in \cite{elkind}, the supported distributions are
    the Gaussian, the Uniform Rectangle, and the Uniform Disc.

  \item \textbf{Survey Data.}
    Upon selecting the "Survey" tab in the sidebar, the user selects two political issue topics to serve as the axes 
    of the plot. Survey responses are plotted using a weighted average and move through the space as users re-rank the survey questions or change the axes topics.
    Candidates that map to the survey questions can either be added randomly (responses to all questions are randomized), or through the user manually answering each
    survey quesiton themselves to create a 'Profile candidate'.
    Details on the weighted average and topic selection are described in Section \ref{sub:Survey Data}.

  \item \textbf{Data Upload.}
    If a user has exported a JSON file of the candidate and voter points,
    or survey respondent points, from a previously-generated plot,
    they may upload that data.

\end{itemize}


\subsection{Survey Data}%
\label{sub:Survey Data}
The Fall 2014 Simon Poll \cite{simon}, conducted by the non-partisan Paul Simon Public
Policy Institute, interviewed 1,006 registered voters across Illinois from
September 23--October 15, 2014.
%
The survey asks respondent for demographic details---such as
sex, race, income level, education level; their plans for voting in the upcoming election;
and for their opinions on political issues including
government reform, LGBTQ+ rights, abortion rights, gun control laws, etc.

The survey code book lists survey questions under topics, such as 'Illinois Political Reform' and 'Guns and Society'.
We use these as a starting point to develop our axis topics. Each topic cluster given by the codebook
underwent pairwise cross-tabulation analysis: for each pair of questions within a topic,
we computed the conditional probability that a respondent holding a conservative
(or liberal) position on one question also held a consistent position on the
other. A question was included only if one political side of pairwise conditional probabilities
exceeded 50\%. This led us to the three political topics we use as axes:
"Role of Government", "Social Issues", and "Guns".


The questions in the survey were originally on
4-point Likert scale (strongly favor to strongly oppose), 3-point scale, and binary scale.
To normalize, we recoded all scales to a 0-1 range before scoring; 
(0, 0.33, 0.67, 1), (0, 0.5, 1), (0,1) respectively.
We use rank-sum weighting based on the user's ranking of the survey questions within each topic cluster.
The weighted average of a respondent's answers to the $k$ questions in the 
topic cluster determines their position on the corresponding axis, with the position changing as
users re-rank survey questions.

\subsection{Implementation}%
\label{sub:Implementation}
The application front end is built with Svelte and D3 in JS,
with Vite as the server.
The election algorithms are implemented in Python and
compiled to Web Assembly via Pyodide so that they can be run in browser.

